%% =====================================================================
%%  T-25-01  The Reneged Price
%%  -------------------------------------------------------------------
%%  One idea per file. Core is mandatory. Slots in canonical order.
%%  Max 700 words. No layout commands. No filler. New source material
%%  for this entry goes in sources/<BOOK>/T-25-01.tex -- never by
%%  rewriting this file (docs/RULES.md R0.1).
%% =====================================================================
%% @kind: topic
%% @id: T-25-01
%% @title: The Reneged Price
%% @subtitle: The terms change after the decision is made, and the decision survives
%% @chapter: c25
%% @order: 10
%% @status: stable
%% @sources: B4
%% @updated: 2026-09-19

\topic{T-25-01}{The Reneged Price}{The terms change after the decision is made, and the decision survives}

\begin{core}
Secure the decision on attractive terms, then reveal that the terms were wrong. The decision frequently stands anyway, on terms nobody would have accepted at the outset. The gap between the offer that produced the yes and the offer that is actually delivered is the profit.
\end{core}
\begin{mechanism}
The technique works because a decision, once made, is not continuously re-evaluated against its original grounds. Commitment supplies its own support: the person has already begun assembling reasons, imagining the outcome and describing themselves as someone who does this thing. Those reasons are attached to the decision rather than to the price, so when the price moves the reasons stay.

Two conditions make it reliable. The change must arrive after the decision and not before it, so that it presents as an unfortunate discovery rather than as the actual offer. And it must be small enough to fall inside the range the person can absorb without reopening the whole question --- which is a range set by how much they have already invested, and grows with it.
\end{mechanism}
\begin{tells}
  \item The terms deteriorate after you have agreed and not before.
  \item The deterioration is presented as an error, an oversight or an unavoidable change rather than as a revision.
  \item You find yourself accepting terms you would have refused at the outset.
  \item The change is small relative to what you have already invested.
  \item Nobody offers to undo the decision along with the terms.
\end{tells}
\begin{moves}
  \item Present the attractive terms first and let the decision be taken on them.
  \item Reveal the correction after signature, deposit or delivery, not before.
  \item Keep the revision inside the range the counterpart can absorb without restarting, and frame it as a discovery you both share.
  \item Offer a small concession alongside the correction so that the exchange still looks reciprocal.
\end{moves}
\begin{counters}
  \item Make the decision conditional on the terms in writing, so that a change in terms is visibly a change in the decision's grounds.
  \item When the price moves, restart the decision rather than adjusting it. Ask whether you would buy this, at this, having not already bought anything.
  \item Pay nothing before the terms are fixed, and treat any deposit as the point at which your own judgement starts being used against you.
  \item Name the pattern plainly. It loses most of its force the moment both parties agree that is what is happening.
  \item Compare against the market rather than against the original offer, which is a comparison you no longer control.
\end{counters}
\begin{cost}
This is the entry in the book most likely to be actionable rather than merely unpleasant: a material term changed after agreement is in many jurisdictions misrepresentation, and the written record that makes the technique work is the same record that proves it. Commercially it also burns the customer --- the transaction succeeds and the relationship does not, which is a bad trade in any market with repeat business or referrals.
\end{cost}

\sources{B4}
\seesources
\seealso{T-14-03, T-14-04, T-24-02, T-19-02}
